\documentclass[12pt]{article}
\usepackage[T1]{fontenc}
\usepackage[utf8]{inputenc}
\usepackage{lmodern}
\usepackage{textcomp}
\usepackage{enumitem}
\usepackage{geometry}
\geometry{margin=1in}
\begin{document}
\begin{center}
Answer Key - Constitutional Law - Graduate Level Assessment 3 \
\small (Answers are ordered exactly as in the question paper.)
\end{center}

\section*{Multiple Choice}
\begin{enumerate}[label=\arabic*.]
\item \textbf{Answer:} E
\\ \textit{Options:} A) The UK Constitution is based on a Bill of Rights; B) The UK Constitution gives the judiciary the power to overturn acts of parliament; C) The UK Constitution's power is derived solely from Parliament; D) The UK Constitution is codified and contained in a single document; E) The UK Constitution is uncodified and can be found in a number of sources
\\ \textit{Note:} Correct option: E - The UK Constitution is uncodified and can be found in a number of sources
\item \textbf{Answer:} A
\\ \textit{Options:} A) Pound; B) Austin; C) Savigny; D) Hart; E) Bentham
\\ \textit{Note:} Correct option: A - Pound
\item \textbf{Answer:} C
\\ \textit{Options:} A) Standard-setting means setting certain standards of conduct in human rights treaties; B) Standard-setting means putting forward both binding and non-binding legal standards; C) Standard-setting means putting forward non-binding legal standards; D) Standard-setting means enforcing human rights treaties; E) Standard-setting means establishing a set of universal human rights principles
\\ \textit{Note:} Correct option: C - Standard-setting means putting forward non-binding legal standards
\item \textbf{Answer:} D
\\ \textit{Options:} A) The Court held that the right to a fair trial trumped the privilege of immunity; B) The Court held that the privilege of immunity was not applicable in this case; C) The Court held that the right to a fair trial was not applicable in this case; D) The Court held that immunities were not in conflict with the right to a fair trial; E) The Court held that the case was admissible due to overriding human rights considerations
\\ \textit{Note:} Correct option: D - The Court held that immunities were not in conflict with the right to a fair trial
\item \textbf{Answer:} A
\\ \textit{Options:} A) Vengeance.; B) Deterrence.; C) Restoration.; D) Moral education.; E) Indemnification.
\\ \textit{Note:} Correct option: A - Vengeance.
\end{enumerate}

\section*{True / False}
\begin{enumerate}[label=\arabic*.]
\setcounter{enumi}{5}
\item \textbf{Answer:} False
\\ \textit{Options:} A) True; B) False
\\ \textit{Note:} LLM-generated false statement. Correct answer: 'Standard-setting means putting forward non-binding legal standards'.
\item \textbf{Answer:} True
\\ \textit{Options:} A) True; B) False
\\ \textit{Note:} LLM-generated true statement based on MCQ.
\item \textbf{Answer:} True
\\ \textit{Options:} A) True; B) False
\\ \textit{Note:} LLM-generated true statement based on MCQ.
\item \textbf{Answer:} False
\\ \textit{Options:} A) True; B) False
\\ \textit{Note:} LLM-generated false statement. Correct answer: 'Goodwin v UK (2002)'.
\item \textbf{Answer:} False
\\ \textit{Options:} A) True; B) False
\\ \textit{Note:} LLM-generated false statement. Correct answer: 'Law is determinate.'.
\end{enumerate}

\section*{Long Form Response}
\begin{enumerate}[label=\arabic*.]
\setcounter{enumi}{10}
\item \textbf{Answer:} The ICJ was not hostile to the accumulation theory. Explain why this is the correct answer and provide supporting evidence. Reference key facts or concepts that support this choice.
\\ \textit{Note:} Long-form derived from MCQ; award credit for stating the correct idea and giving supporting reasoning.
\item \textbf{Answer:} admitted, because it is relevant to the case.. Explain why this is the correct answer and provide supporting evidence. Reference key facts or concepts that support this choice.
\\ \textit{Note:} Long-form derived from MCQ; award credit for stating the correct idea and giving supporting reasoning.
\end{enumerate}

\vfill
\noindent\textit{End of Answer Key}
\end{document}